The conclusion was weakened on 2026-10-23: Corollary \claimref{standard-model.one-vev-vector-relations-breaking-pattern} reads the up-type masses through the conjugate $c^{*}$, and through $c^{*}$ one rung vector gives $m_{u} = m_{d}$ whether it is real or complex, so the computation below, which gives $b = \bar{a}$ for a real vector, no longer shows that the rung field must be complex; what follows from $m_{t} \neq m_{b}$ is a second independent doublet coupling, which a complex bidoublet, carrying two electroweak doublets, supplies. Here the imaginary unit is $I_{4}$, which commutes with every mode operator, so $\gamma_{7} = I_{4}(v_{4} - w_{4})$ and $x(w_{4} + v_{4}) + y\,I_{4}(v_{4} - w_{4}) = (x + yI_{4})\,v_{4} + (x - yI_{4})\,w_{4}$, the two coefficients being conjugate in $\mathrm{span}\{1, I_{4}\}$, so $|a| = |b|$ and $m_{u} = m_{d}$ at tree level by \claimref{standard-model.one-vev-vector-relations-breaking-pattern}; the script checks it with $x = 3$, $y = 5$. The argument holds given \claimref{standard-model.hyp-s-slot-statistics}, under which the sectors of the module are the species of the generation, so that the two magnitudes are the masses $m_{u}$ and $m_{d}$.
